\originalsection{磁通变化的三种来源与开路棒}\label{tb:motional}
均匀磁场穿过平面线圈时, $\Phi=BA\cos\theta$, 所以
\[\mathcal E=-A\cos\theta\dot B-B\cos\theta\dot A+BA\sin\theta\dot\theta.\]
三个项依次对应变场、变面积、转动. 只有每匝通量都相同才将整个式乘$N$; 通量不同则先求磁通链之和. 先规定法向, 再求带符号的$\Phi$和$\dot\Phi$, 最后由$\mathcal E=-\dot\Phi$判环向. 磁通为负但变得更负时, 感应环向与负磁通变小的情况相反.

孤立导棒以$\vct v=v\uv x$运动, $\bb=B\uv z$, 正载流子受磁偏转向$-y$, 电子偏向$+y$, 两端分离出电荷. 开路稳态满足$\ee+\vct v\times\bb=0$, 所以棒内$\ee=vB\uv y$, $+y$端电势较低, 长$\ell$的端势差大小$vB\ell$. 没有闭路不代表没有电荷分离; 闭路时分离场、欧姆电流和局部运动驱动共同满足实际回路. 电动势是沿完整路径的单位电荷驱动功, 与开路端电压在给定测量条件下相关, 两者不应无条件等同.

低速换参考系时, 原来静止场对运动棒的$\vct v\times\bb$驱动在棒静止系可表现为电场. 近似$\ee'\simeq\ee+\vct v\times\bb$. 所以动生和感生的拆分依赖观察系, 同一实验的可观测结果相容. 这里是低速近似, 没有声称不同惯性系的三维力和所有场分量精确相同.

\originalsection{感应电场可以伸到无磁场处}\label{tb:induced-e}
磁场$B(t)\uv z$只在半径$a$的长圆柱区域内近似均匀, 外部磁场在理想模型中为零. 轴对称使感应电场沿$\uv\phi$. 用半径$s$圆周作固定回路:
\[E_\phi(s)=\begin{cases}-s\dot B/2,&s<a,\\ -a^2\dot B/(2s),&s>a.\end{cases}\]
内圈包围面积$\pi s^2$, 外圈只包围有效磁通$\pi a^2B$. $s=a$两侧连续, 外部即使局部$B=0$仍有非零感应$E$. 这个场不靠测试导线才存在; 放上环只是测出环量. 它的闭路积分不零, 所以不能在包围变化磁通的区域定义单值静电势来替代整条路径的驱动. 原先电势差的路径独立条件在这里已经失效.

\begin{example}
\textbf{原10.9.5: 圆环中的线性变场.} 
半径$a$的环垂直于指入纸面的$B(t)=B_0+bt$, $B_0,b>0$, 总电阻$R$, 忽略自感. 求$t=0$通量、电动势、电流方向及热功率.
\end{example}
\begin{solution}
法向取入纸面、环向相应顺时针, $\Phi(0)=\pi a^2B_0$, $\mathcal E=-\pi a^2b$. 负号表示逆时针电流, 大小$I=\pi a^2b/R$, 热功率$P=\pi^2a^4b^2/R$. 增强的入纸面场被感生的出纸面场阻碍. 若自感不可忽略, 须先解$L\dot I+RI=\mathcal E$而非直接除$R$.
\end{solution}

\originalsection{长直线旁的磁通: 感应与互感共用一次积分}\label{tb:wire-flux}
\begin{example}
\textbf{原10.9.1、10.9.4、10.9.6及11.11.3.} 
原图长直线电流在纸面向上、矩形环在导线右侧, 环处原场指入纸面. 环长$\ell$沿导线方向, 宽$w$, 近边距离$s$. 取法向为导线在环处场的方向. 分别求: (a)导线电流$I_1(t)=a+bt$、环静止时的感应电动势; (b)平行于导线、长度$\ell$的开路棒以速度$v$远离导线时的端电动势; (c)整个矩形环以$v=\dot s$远离稳恒电流线时的感应电流; (d)固定几何的互感.
\end{example}
\begin{solution}
先对非均匀场积分, $\dd A=\ell\dd r$:
\[\Phi=\int_s^{s+w}\frac{\mu_0I_1\ell}{2\pi r}\dd r
=\frac{\mu_0I_1\ell}{2\pi}\ln\frac{s+w}s.\]
(a)$\mathcal E=-\mu_0b\ell\ln[(s+w)/s]/(2\pi)$, 环流与增大的原磁通反向, 从纸面看为逆时针. (b)同一瞬间沿平行棒的轴距相同, 故$|\mathcal E|=\mu_0I_1\ell v/(2\pi s)$, 极性按$\vct v\times\bb$确定. 原10.9.4图将导线画成水平、电流向左, 棒位于导线上方并向上远离; 场入纸面, 正电荷受力向左, 所以左端正、右端负.

(c)电流线稳恒, 环的外移使磁通减小:
\[\mathcal E=\frac{\mu_0I_1\ell v}{2\pi}\left(\frac1s-\frac1{s+w}\right)
=\frac{\mu_0I_1\ell wv}{2\pi s(s+w)}.\]
环电流大小为该量除$R$, 方向保持入纸面的原磁通, 从纸面看为顺时针. 也可将近远两竖边的动生电动势相减; 横边$\dd\vct\ell$与$\vct v\times\bb$垂直而无贡献. 两法得到同一式, 这给移动边界通量公式的直接检验.

(d)$M=\Phi/I_1=\mu_0\ell\ln[(s+w)/s]/(2\pi)$. 绕向反转使带符号$M$反号. 前一个感应题中的几何系数, 就是固定线圈的互感.
\end{solution}

\begin{example}
\textbf{原10.9.2、10.9.3: 面积改变与共享导棒.} 
(a)四根长度$\ell$的导线构成正方形, 在指入纸面的垂直磁场中变成内角$\theta$的菱形, 历时$\Delta t$. 求平均电动势及电阻$R$中平均电流. (b)一根长度$\ell$导棒在平行导轨上向左以$v$匀速移动, 左右端各有负载$R_1,R_2$, 场指入纸面. 求两电阻电流、热功率和所需外力.
\end{example}
\begin{solution}
(a)面积从$\ell^2$变成$\ell^2\sin\theta$, 通量减少, 感生场保持原方向, 所以电流顺时针. 平均电动势大小$B\ell^2(1-\sin\theta)/\Delta t$, 平均电流大小为它除$R$. 若题给的是瞬时$\theta(t)$, 则电动势为$-B\ell^2\cos\theta\dot\theta$, 不能把总变化除历时当每一瞬间的值.

(b)导棒相当于两支路共用的动生源, 幅度$\mathcal E_0=B\ell v$. 左回路顺时针、右回路逆时针, 但共享棒中的电流同向相加:
\[|I_1|=\frac{B\ell v}{R_1},\quad |I_2|=\frac{B\ell v}{R_2},\quad
P=B^2\ell^2v^2\left(\frac1{R_1}+\frac1{R_2}\right).\]
磁力向右, 大小$B^2\ell^2v(1/R_1+1/R_2)$, 外力须向左且同大, 外力功率恰为$P$. 共享杆电流不能按两个回路的反向环流相减.
\end{solution}

\originalsection{磁阻尼、发电机及涡流的动力学}\label{tb:magnetic-braking}
撤去匀速滑杆的外力, 单回路电阻$R$且忽略自感时,
\[m\dot v=-\frac{B^2\ell^2}R v,\quad v=v_0e^{-t/\tau_m},\quad
\tau_m=\frac{mR}{B^2\ell^2}.\]
位移为$v_0\tau_m(1-e^{-t/\tau_m})$, 无限时间总滑行距离有限. 若惯性过程中离开磁场区, 指数模型只能用到离开为止. 热量积分$\int_0^\infty B^2\ell^2v^2\dd t/R=mv_0^2/2$, 说明机械能转热而全系统能量守恒.

沿倾角$\theta$的斜轨滑动, 若磁场竖直向下, 正对轨道面的分量是$B\cos\theta$. 向下速度$v$时制动力$B^2\ell^2\cos^2\theta\,v/R$, 所以终速
\[v_t=\frac{mgR\sin\theta}{B^2\ell^2\cos^2\theta}.\]
达到终速时重力功率$mgv_t\sin\theta$等于热功率. 若电池驱动水平滑杆, 反电动势$B\ell v$抵消电池电动势的一部分, $I=(\mathcal E-B\ell v)/R$, 从静止
\[v(t)=\frac{\mathcal E}{B\ell}(1-e^{-t/\tau_m}).\]
这两个模型的驱动力不同, 不能沿用同一个终速. 回路自感显著时, 要联立电流和速度的微分方程.

发电机线圈$N$匝、面积$A$、恒速$\omega$, 通量链$NBA\cos\omega t$, 纯电阻负载下$I=NBA\omega\sin\omega t/R$. 电流的磁矩使转动力矩阻碍转动, 外力矩的瞬时功率
\[P_{\rm mech}=\mathcal EI=\frac{(NBA\omega)^2}{R}\sin^2\omega t=I^2R.\]
固定转速下平均功率为$(NBA\omega)^2/(2R)$, 所以负载电阻越小, 取电及所需平均机械功率越大. 理想无负载时只有开路电压, 并没有这个负载热功率.

实体金属内可形成涡流. 在局部磁场边缘移动或非均匀时间变场中, 许多内部回路的通量改变, 涌现环流和磁阻尼. 开槽切断大回路、叠片间加绝缘可以提高环流电阻和缩小回路面积, 减少铁芯损耗. 整个闭合刚性金属环在无限均匀静磁场中平移, 磁通不变, 两侧动生驱动相消, 并不因为“有运动”就必有净电流.

\originalsection{自感的几何计算与双重能量核对}\label{tb:toroid-l}
\begin{example}
\textbf{原11.3及11.11.1: 环形线圈的自感与能量.} 
理想环形绕组$N$匝, 矩形截面高$h$, 内外半径$a,b$, 真空芯. 求自感及储能, 用场积分核对.
\end{example}
\begin{solution}
内部$B_\phi=\mu_0NI/(2\pi r)$, 穿过每匝的截面磁通是
\[\Phi=\int_a^bB_\phi h\dd r=\frac{\mu_0NIh}{2\pi}\ln\frac ba,
\quad L=\frac{N\Phi}I=\frac{\mu_0N^2h}{2\pi}\ln\frac ba.\]
故$U=\mu_0N^2I^2h\ln(b/a)/(4\pi)$. 另一算法是对圆柱小壳体积$2\pi rh\dd r$积分:
\[U=\int_a^b\frac{B^2}{2\mu_0}2\pi rh\dd r=\frac{\mu_0N^2I^2h}{4\pi}\ln\frac ba.\]
$b-a\ll a$时$\ln(b/a)\simeq(b-a)/a$, 还原长为$2\pi a$、截面$h(b-a)$的长管型式. 不薄时场并不恒定, 不能只取中心$B$乘整个截面. 理想绕组内部之外的场近零是密绕近似, 实际稀疏匝及引线有漏场.
\end{solution}

\begin{example}
\textbf{原11.11.2: 实导线内部的磁能.} 
非磁性圆线半径$R$, 长$\ell$, 稳恒电流$I$在截面均匀分布. 求线内磁能及内部电感.
\end{example}
\begin{solution}
$B(s)=\mu_0Is/(2\pi R^2)$. 因此
\[U_{\rm int}=\int_0^R\frac{B^2}{2\mu_0}2\pi s\ell\dd s
=\frac{\mu_0I^2\ell}{4\pi R^4}\int_0^Rs^3\dd s=\frac{\mu_0I^2\ell}{16\pi}.\]
与$\tfrac12L_{\rm int}I^2$比较得$L_{\rm int}/\ell=\mu_0/(8\pi)$. 高频趋肤效应改变电流分布, 内部能量不再等于该稳恒值. 第\ref{ch:induction}章原有同轴外部电感则只积分$a<s<b$的绝缘区, 两部分须在条件一致时相加.
\end{solution}

\originalsection{互感、耦合与总磁通的保存}\label{tb:coupled-l}
原例11.1同面同心大、小环半径$R_1,R_2$, 仅在$R_2\ll R_1$时小环面内近似有大环中心匀场. 因而$M\simeq\mu_0\pi R_2^2/(2R_1)$. 原例11.4长管长$\ell$, 内绕组$N_1$匝, 有$N_2$匝的外绕组链接同一截面$A$, 得$M=\mu_0N_1N_2A/\ell$. 外绕组只需计实际链接匝数, 松绕和漏磁不能直接当完全同磁通的密绕模型.

两线圈按一致端点符号, 总能量
\[U=\tfrac12L_1I_1^2+\tfrac12L_2I_2^2+MI_1I_2.\]
被动系统任意电流下能量非负, 完全平方或二次式判别式给$M^2\le L_1L_2$. 耦合系数$k_c=M/\sqrt{L_1L_2}$保留带符号, $|k_c|\le1$. 串联同向电流给$L_{\rm eq}=L_1+L_2+2M$, 反向连接给减号. 无互感并联时电压相同,$L_1\dot I_1=L_2\dot I_2$, 所以$1/L_{\rm eq}=1/L_1+1/L_2$. 有互感并联则须先解
\[V=L_1\dot I_1+M\dot I_2=L_2\dot I_2+M\dot I_1,\]
在$L_1+L_2-2M\ne0$且对应连接约定下得$L_{\rm eq}=(L_1L_2-M^2)/(L_1+L_2-2M)$; 理想退化条件须另外处理.

导电环受到外磁通$\Phi_{\rm ext}(t)$时, 完整慢变回路式是
\[L\dot I+RI=-\dot\Phi_{\rm ext}.\]
$R=0$的理想闭合环保持$LI+\Phi_{\rm ext}$这个\emph{总磁通链}恒定, 并非让外场停止改变. 有限$R$时环流衰减并产生热, 磁体或电源还受反作用力. 感生电流结果若使$\omega L$与$R$相当, 原先忽略自感的除$R$算法就需返修, 这是模型自检而非无限缩小电阻就可获得无限电流.

\originalsection{磁压力与张力的可测含义}\label{tb:magnetic-pressure}
两平行无限电流面大小$K$, 方向相反, 场在中间大小$B=\mu_0K$, 外部相消. 每张面对自身产生的场不能取作自力, 另一面在它处产生大小$\mu_0K/2$. 因而两面相斥, 压力$K(\mu_0K/2)=B^2/(2\mu_0)=u_B$. 同向电流面则相吸. 磁场表现出沿场张力和横向压力, 原11.9中的线圈匝间力、长管向外扩张可以从它理解. 均匀场的局部应力存在, 一个无电流小区域闭面上各边应力积分仍可相消, 不能仅凭$u_B>0$就推出区域合力非零.

原11.4(b)在两组绕组理想完全耦合、链接同一均匀磁路且有效长度为$\ell$时,
\[L_1=\frac{\mu_0N_1^2A}{\ell},\quad L_2=\frac{\mu_0N_2^2A}{\ell},\quad
M=\frac{\mu_0N_1N_2A}{\ell}=\sqrt{L_1L_2}.\]
取同向绕组和正互感得到正根, 反绕时$M=-\sqrt{L_1L_2}$. 这完成该例的自感/互感关系, 一般漏磁时只保留之前的界限, 不自动等号.

若一根直棒长$\ell$绕固定端在垂直匀场$B\uv z$中以$\omega$转动, 各段速度为$\omega r$, 不能整棒使用端点速度. 径向动生驱动积分
\[\mathcal E_{\rm center\to end}=\int_0^\ell(\omega r\uv\phi\times B\uv z)\cdot\uv r\dd r=\tfrac12B\omega\ell^2.\]
正电荷被推向外端, 外端高势. 这建立原10.11.5转棒方法, 也适用于转盘沿半径的开路势差; 闭合回路仍须把其余运动与静止部分都算进去.
